% SPDX-License-Identifier: GPL-3.0-or-later
% Copyright (C) 2026 Alessandro Greco
\chapter{«Луковица»: шифрование ячеек RELAY и CGO}\label{cap:10-cifratura-relay}
\aggiornatoal{23 сентября 2026 г.}{коммит torspec \texttt{8e1f245}; tor 0.4.9.13; предложение 359}

\begin{obiettivi}
К концу этой главы вы сможете:
\begin{itemize}[nosep]
  \item описывать, как клиент оборачивает сообщение одним слоем шифрования для каждого перехода и как реле снимают эти слои;
  \item объяснять роль Recognized и накопительного дайджеста и почему нужны оба механизма;
  \item определять слабости традиционной схемы: пластичность и атаки \emph{маркирования (tagging)};
  \item описывать базовый принцип Counter Galois Onion (CGO), схемы, которая её заменяет;
  \item с помощью собственного кода расширять реальную цепь до трёх переходов (лабораторная работа~4).
\end{itemize}
\end{obiettivi}

В \cref{cap:09-circuiti} клиент согласовал с каждым реле цепи четыре ключа: Kf и Kb для шифрования, Df и Db для проверки целостности. Теперь задействуем их. Именно этот механизм дал onion routing его название: сообщение оборачивается в столько слоёв, сколько реле входит в маршрут, а каждое реле снимает один слой, словно очищая луковицу.

% ---------------------------------------------------------------------------
\section{Традиционная схема}\label{sec:tor1}

Схема, используемая с первых версий Tor и называемая в новом коде \emph{tor1}, сочетает для каждого перехода и каждого направления два средства:
\begin{itemize}
  \item потоковый шифр --- \textbf{AES-128 в режиме счётчика} (CTR), с ключом Kf или Kb;
  \item \textbf{накопительный дайджест SHA-1}, инициализированный Df или Db.
\end{itemize}
Оба механизма сохраняют \emph{состояние} в течение всей жизни цепи. Счётчик AES-CTR не начинается с нуля для каждой ячейки: первая использует первые 509 байт гаммы шифрования, вторая --- следующие 509 и так далее. Аналогично дайджест вычисляется не для одной ячейки, а для \emph{всех} ячеек, которыми до этого обменивались с данным переходом. Поэтому клиент и реле обязаны обрабатывать ячейки строго в одном порядке: потеря или дублирование ячейки рассинхронизирует оба состояния. TLS, гарантирующий порядок и целостность на канале, делает это возможным.

\subsection{Прямое направление: от клиента внутрь цепи}

Чтобы отправить сообщение переходу $d$, клиент~\spec{tor-spec}{Routing relay cells}:
\begin{enumerate}
  \item готовит открытое тело сообщения RELAY (\cref{sec:relay}) с Recognized и Digest, равными нулю;
  \item обновляет дайджест прямого направления перехода $d$ этим телом и записывает первые четыре байта результата в поле Digest;
  \item шифрует тело ключом Kf перехода $d$, затем Kf перехода $d-1$ и так далее до перехода 1: внешний слой принадлежит первому реле.
\end{enumerate}
Каждое реле, получив ячейку, расшифровывает её своим Kf, снимая один слой, и проверяет, предназначена ли ячейка ему. Если нет, оно пересылает ячейку дальше, где следующий переход снимает ещё один слой. Слои показаны на \cref{fig:cipolla}.

\begin{figure}[htbp]
\centering
% SPDX-License-Identifier: GPL-3.0-or-later
% Copyright (C) 2026 Alessandro Greco
%% Русское альтернативное описание EPUB.
\altfig{Шифрование RELAY в прямом направлении (tor1): клиент применяет Kf3, Kf2, Kf1, а сторожевой, промежуточный и выходной узлы снимают соответствующие слои. Внутренние блоки всё ещё содержат шифротекст. Выходной узел распознаёт RELAY при нулевом Recognized и верном накопленном дайджесте. Данные приложения могут оставаться зашифрованными HTTPS. Метаданные трафика на схеме не показаны.}
% Слои шифрования в прямом направлении.
\begin{tikzpicture}[font=\sffamily\small,
  s1/.style={draw=torgrigio,fill=torgrigio!15,rounded corners=3pt},
  s2/.style={draw=torviola!70,fill=torviola!18,rounded corners=3pt},
  s3/.style={draw=torverde!60!black,fill=torverde!28,rounded corners=3pt},
  msg/.style={font=\ttfamily\tiny,fill=white,draw,inner sep=2pt}]
  % Пакет с несколькими слоями
  \newcommand{\pacchetto}[2]{%
    \ifnum#2>2 \draw[s1] (#1-1.05,-0.5) rectangle (#1+1.05,0.5);\fi
    \ifnum#2>1 \draw[s2] (#1-0.85,-0.37) rectangle (#1+0.85,0.37);\fi
    \ifnum#2>0 \draw[s3] (#1-0.65,-0.25) rectangle (#1+0.65,0.25);\fi
    \ifnum#2=0 \node[msg] at (#1,0) {сообщение};\else \node[msg] at (#1,0) {шифротекст};\fi}
  \pacchetto{0}{3}
  \node[relay] (g) at (2.45,0) {Сторожевой};
  \pacchetto{4.35}{2}
  \node[relay] (mi) at (6.2,0) {Промежуточный};
  \pacchetto{7.9}{1}
  \node[relay,fill=torverde!30] (e) at (9.6,0) {Выходной};
  \draw[cellaf] (1.05,0) -- (g); \draw[cellaf] (g) -- (3.5,0);
  \draw[cellaf] (5.2,0) -- (mi); \draw[cellaf] (mi) -- (7.25,0);
  \draw[cellaf] (8.55,0) -- (e);
  % После выходного узла
  \begin{scope}[yshift=-1.6cm]\pacchetto{9.6}{0}\end{scope}
  \draw[cellaf] (e) -- (9.6,-1.4);
  \node[font=\scriptsize,align=center] at (0,-0.95) {клиент: шифрует Kf$_3$,\\затем Kf$_2$, затем Kf$_1$};
  \node[font=\scriptsize] at (2.45,-0.85) {снимает Kf$_1$};
  \node[font=\scriptsize] at (6.2,-0.85) {снимает Kf$_2$};
  \node[font=\scriptsize,align=right,anchor=east] at (8.8,-1.6) {снимает Kf$_3$: Recognized = 0\\и дайджест совпадает};
  % Легенда
  \draw[s1] (1.2,1.0) rectangle ++(0.3,0.2); \node[font=\scriptsize,anchor=west] at (1.55,1.1) {слой сторожевого};
  \draw[s2] (3.9,1.0) rectangle ++(0.3,0.2); \node[font=\scriptsize,anchor=west] at (4.25,1.1) {слой промежуточного};
  \draw[s3] (6.0,1.0) rectangle ++(0.3,0.2); \node[font=\scriptsize,anchor=west] at (6.35,1.1) {слой выходного};
\end{tikzpicture}

\caption{Шифрование RELAY по традиционной схеме tor1: клиент применяет Kf$_3$, Kf$_2$, Kf$_1$, а каждый ретранслятор снимает свой слой. Выходной узел распознаёт сообщение при Recognized = 0 и правильном накопленном дайджесте; данные приложения могут оставаться зашифрованными HTTPS. Вложенные рамки условно изображают шифротекст.}
\label{fig:cipolla}
\end{figure}

\subsection{Обратное направление: от цепи к клиенту}

В обратном направлении роли меняются. Реле, порождающее сообщение, снабжает его дайджестом обратного направления и шифрует ключом Kb. Каждое следующее реле на пути к клиенту \emph{добавляет} ещё один слой с собственным Kb. Клиент, знающий все ключи, снимает слои по одному, начиная с первого перехода, и после каждого слоя проверяет, не этим ли переходом была создана ячейка.

\subsection{Цепь как \enquote{труба с промежуточными выходами (leaky pipe)}}

Клиент может отправить сообщение \emph{любому} переходу цепи, а не только последнему, просто остановив шифрование на нужном переходе. Именно так сообщение \texttt{EXTEND2} для третьего перехода доставляется второму реле. Спецификация называет это \emph{leaky pipe topology}: цепь подобна трубе, из которой сообщения могут выйти до конца. Исходная статья Tor~\cite{dingledine2004tor} предлагала использовать это свойство и как защиту: отправка некоторых потоков через промежуточное реле затруднила бы наблюдателю на выходе определение соответствующей цепи. На практике tor выводит потоки только через последний переход.

% ---------------------------------------------------------------------------
\section{Распознавание ячейки}\label{sec:riconoscere}

Как реле после снятия своего слоя узнаёт, адресована ли ячейка ему? Нельзя просто \enquote{проверить, имеет ли содержимое смысл}: для программы и всё ещё зашифрованная, и открытая ячейка являются последовательностями байтов. Традиционная схема выполняет две проверки подряд~\spec{tor-spec}{Relay cells}:

\begin{enumerate}
  \item \textbf{Recognized}: эти два байта должны быть нулевыми. Если ячейка всё ещё зашифрована, байты псевдослучайны и равны нулю лишь один раз на \num{65536}. Проверка почти бесплатна и сразу отбрасывает подавляющее большинство транзитных ячеек.
  \item \textbf{Digest}: если Recognized равен нулю, реле вычисляет накопительный дайджест \emph{так, как если бы} ячейка принадлежала ему, при этом поле Digest обнуляется, и сравнивает первые четыре байта с содержимым поля. Если они совпадают, ячейка \enquote{распознана}: реле обновляет состояние дайджеста и обрабатывает ячейку.
\end{enumerate}

Деталь, которую легко упустить: если Recognized равен нулю, но дайджест не совпадает, это ложное срабатывание, и состояние дайджеста \textbf{не должно} обновляться. В лабораторной работе~4 мы вычисляем на копии состояния (\texttt{self.digest.copy()}) и сохраняем её только при успехе.

Дайджест охватывает всё тело, включая байты заполнения, которые поэтому являются случайными (\cref{sec:relay}). Для цепей onion-сервисов, где клиент и сервис делят ключи \enquote{виртуального перехода} (\cref{cap:14-onion-services}), вместо SHA-1 используется SHA3-256.

Что происходит, если ячейку никто не распознал? Если последнее реле получает прямую ячейку и не распознаёт её, произошла ошибка: оно уничтожает цепь посредством \cella{DESTROY}. Если ни один переход не распознаёт входящую к клиенту ячейку, клиент делает то же самое.

% ---------------------------------------------------------------------------
\section{Слабости tor1}\label{sec:tor1-debolezze}

\subsection{Пластичность}

У режима счётчика есть неудобное свойство: шифротекст равен $C = P \oplus S$, где $S$ --- гамма шифрования. Изменение одного бита $C$ меняет тот же бит $P$, не требуя знания ключа. Если противник знает или может угадать часть открытого текста, он способен заменить её выбранным текстом: достаточно XOR-нуть шифротекст с разностью двух текстов. В лабораторной работе~4 есть тест, который делает именно это, меняя \enquote{заплати 10 евро} на \enquote{заплати 99 евро}.

В tor1 целостность проверяется только \emph{в пункте назначения}: целевое реле проверяет дайджест, а промежуточные реле этого не делают. Изменённая ячейка поэтому может беспрепятственно пройти до назначения, где дайджест перестанет совпадать.

\subsection{Атака маркирования (tagging)}

Вместе эти свойства допускают специфическую атаку \emph{маркирования (tagging)}. Представим противника, контролирующего входное и выходное реле цепи:
\begin{enumerate}
  \item Входное реле меняет несколько битов транзитной ячейки по узнаваемому шаблону --- \enquote{метке}, или \emph{tag}.
  \item Ячейка проходит через промежуточное реле, которое не может обнаружить изменение.
  \item До проверки дайджеста выходное реле проверяет наличие метки. Если она есть, выходное реле узнаёт, что цепь проходит через сотрудничающий входной узел, и может связать пункт назначения трафика с адресом клиента, известным этому входному узлу.
\end{enumerate}
Если выходное реле не участвует в сговоре, дайджест не совпадёт, цепь будет уничтожена, а клиент построит новую. Противник может повторить попытку, и каждая неудача что-то сообщает. Метку можно накладывать и в обратном направлении, от выходного реле к клиенту.

\begin{approfondimento}[Насколько важно маркирование?]
Противник, контролирующий одновременно входное и выходное реле, уже может коррелировать временные характеристики и объёмы трафика без метки (\cref{cap:26-attacchi}). Поэтому разработчики Tor долго считали маркирование реальной, но менее приоритетной проблемой~\cite{dingledine2004tor}. Однако маркирование надёжнее статистической корреляции и работает даже при очень малом объёме трафика. Предложения по его устранению появлялись более десяти лет: предложения 202, 261 (на основе шифра AEZ), 295 и 308, а затем 359, приведшее к CGO.
\end{approfondimento}

\subsection{Другие ограничения}

Дайджест имеет длину лишь четыре байта, а лежащий в его основе SHA-1 больше не считается безопасным для новых применений. Кроме того, схема не обеспечивает прямую секретность \emph{внутри} цепи: тот, кто получит ключи перехода, пока цепь активна, сможет расшифровать и все ранее записанные ячейки, поскольку гамма шифрования полностью определяется начальным ключом.

% ---------------------------------------------------------------------------
\section{Counter Galois Onion}\label{sec:cgo}

\textbf{CGO} (\emph{Counter Galois Onion}) --- схема шифрования ячеек RELAY, разработанная Дегабриэле, Меллони, Мюнхом и Стамом и принятая в \proposta[Предложении]{359}. Здесь даётся функциональный обзор; полная криптографическая конструкция описана в предложении и работах авторов.

\subsection{Свойства}

Согласно предложению, CGO даёт три улучшения:
\begin{itemize}
  \item схема \textbf{непластична}: любое изменение ячейки делает эту и все последующие ячейки цепи нерасшифровываемыми. Метку нельзя ни прочитать, ни удалить: она способна только разрушить цепь;
  \item обеспечивается \textbf{прямая секретность на уровне ячеек}: после обработки ячейки у стороны больше нет материала, необходимого для её повторного расшифрования;
  \item аутентификация имеет стойкость 128 бит вместо значительно более короткой проверки традиционной схемы.
\end{itemize}
Поскольку CGO больше не использует SHA-1, в целом она также должна быть быстрее.

\subsection{Как это работает на высоком уровне}

Базовые компоненты --- AES-128 и POLYVAL, универсальная хеш-функция из того же семейства, которое используется режимом AES-GCM-SIV. На их основе предложение строит \emph{широкоблочный настраиваемый шифр (wide-block tweakable cipher)}: все 509 байт ячейки обрабатываются как один блок, поэтому изменение хотя бы одного бита в любом месте непредсказуемо меняет весь результат. Расшифрованная ячейка состоит из \emph{тега} (16 байт) и тела (493 байта), как мы видели в \cref{sec:nuovo-formato}.

Для каждого перехода и каждого направления клиент и реле разделяют:
\begin{itemize}
  \item 64-байтовый ключ $K$ и 16-байтовое значение $N$, полученные при согласовании ключей цепи;
  \item \emph{тег} $T'$ последней обработанной ячейки, изначально нулевой.
\end{itemize}
Три идеи делают схему устойчивой к маркированию:
\begin{description}
  \item[Распознавание по тегу] После снятия реле своего слоя ячейка принадлежит ему, если полученный тег совпадает с $N$. Это 16-байтовая проверка вместо двух байт Recognized и четырёх байт дайджеста.
  \item[Связывание] Тег предыдущей ячейки входит в обработку следующей как параметр шифра. Изменённая ячейка создаёт неправильный тег, который, в свою очередь, делает невозможным расшифрование всех последующих ячеек: цепь перестаёт работать.
  \item[Обновление ключей] Каждый раз, когда переход распознаёт или создаёт ячейку, он заменяет $K$ и $N$ новыми значениями, выведенными из предыдущих. Старые значения стираются, обеспечивая прямую секретность на уровне ячеек.
\end{description}
Байт команды ячейки (\cella{RELAY} или \cella{RELAY\_EARLY}) включён в вычисления как \emph{связанные данные (associated data)}: реле не может превратить один тип в другой, не сделав ячейку нерасшифровываемой.

\begin{attenzione}[Направления нельзя менять местами]
Конструкция, лежащая в основе CGO, непластична в одном направлении и пластична в другом. Предложение тщательно выбирает, какое направление использовать клиенту и какое реле, и прямо предупреждает разработчиков не менять их местами. Это показывает, как решение, кажущееся деталью реализации, способно подорвать безопасность всей криптографической схемы.
\end{attenzione}

\subsection{CGO на практике}

\begin{itemize}
  \item \textbf{Ключи}: из потока согласования ключей каждое направление получает 64-байтовый $K$ и 16-байтовый $N$ вместо Df, Db, Kf и Kb.
  \item \textbf{SENDME}: аутентифицированное управление потоком из \proposta[Предложения]{289} использовало накопительный дайджест как доказательство получения данных. В CGO такого дайджеста нет, поэтому SENDME вместо него несёт тег последней подтверждённой ячейки (\cref{cap:11-stream}).
  \item \textbf{Согласование}: клиент использует CGO с реле только если реле объявляет субпротокол \texttt{Relay=6}, а клиент запрашивает его при согласовании ключей ntor-v3 (\cref{sec:ntor3}). В \cref{cap:09-circuiti} мы видели, что реле отклоняет запрос, если он не сопровождается запросом управления перегрузкой.
  \item \textbf{Развёртывание}: tor поддерживает CGO начиная с серии 0.4.9, а в Arti он всегда включён. Одна цепь может содержать переходы с CGO и переходы, всё ещё использующие tor1: выбор делается независимо для каждого перехода.
\end{itemize}

% ---------------------------------------------------------------------------
\section{Лабораторная работа: трёхпереходная цепь}\label{sec:lab-cifratura}

\begin{laboratorio}[4]{Послойное шифрование и расширение цепи}
\textbf{Цель.} Написать \file{labs/ru/код/torlab/encryption.py}, реализующий клиентскую сторону tor1, и использовать его для расширения реальной цепи до выходного реле.\\
\textbf{Предварительные условия.} Лабораторные работы~1--3; работающая лабораторная сеть.\\
\textbf{Оценочное время.} 3 часа.
\end{laboratorio}

\paragraph{Шаг 1: одно направление одного перехода.} В центре лабораторной работы находится небольшой класс, объединяющий AES-CTR и дайджест SHA-1 одного направления:
\begin{lstlisting}[language=Python]
class Direction:
    @classmethod
    def new(cls, key, digest_seed):
        d = hashlib.sha1(); d.update(digest_seed)
        return cls(_aes_ctr(key), d)

    def apply(self, data):             # шифрование и расшифрование: CTR выполняет одну операцию
        return self.cipher.update(data)

    def sign(self, body):                 # записать поле Digest
        self.digest.update(_without_digest(body))
        return body[:5] + self.digest.digest()[:4] + body[9:]

    def recognizes(self, body):
        if body[1:3] != b"\x00\x00":
            return False
        trial = self.digest.copy()
        trial.update(_without_digest(body))
        if trial.digest()[:4] != body[5:9]:
            return False                 # ложное срабатывание: состояние не меняется
        self.digest = trial
        return True
\end{lstlisting}
Объект шифра из \texttt{cryptography} сохраняет счётчик между вызовами: именно такое поведение здесь требуется.

\paragraph{Шаг 2: клиент.} Имея по одному объекту \texttt{Direction} на направление и переход, шифрование и расшифрование превращаются в циклы:
\begin{lstlisting}[language=Python]
def encrypt(self, msg, destination):
    body = self.hops[destination].forward.sign(encode_relay(msg))
    for i in range(destination, -1, -1):
        body = self.hops[i].forward.apply(body)
    return body

def decrypt(self, body):
    for i, hop in enumerate(self.hops):
        body = hop.backward.apply(body)
        if hop.backward.recognizes(body):
            return i, decode_relay(body)
    raise CellError("ячейка не распознана")
\end{lstlisting}

\paragraph{Шаг 3: тесты в памяти и их ограничения.} Файл \file{labs/ru/код/torlab/test/test\_encryption.py} строит имитированную цепь из трёх реле. Он проверяет, что каждый переход распознаёт только свои ячейки, промежуточные реле видят лишь псевдослучайные байты, а изменение обнаруживается в пункте назначения; также демонстрируется пластичность AES-CTR. Но у таких тестов есть важное ограничение: имитированные клиент и реле используют \emph{один и тот же код}. Если, например, мы ошибёмся в порядке обновления дайджеста, обе стороны допустят одинаковую ошибку, и тест всё равно пройдёт. Окончательная проверка --- общение с независимой реализацией.

\paragraph{Шаг 4: расширить реальную цепь.} Устанавливаем первый переход к \texttt{test003r} через ntor-v3, как в лабораторной работе~3, затем расширяемся до \texttt{test004r} и выходного узла \texttt{test006e}. Модуль \file{labs/ru/код/torlab/laboratory.py} читает идентичности и ключи ntor реле из файлов лабораторной сети:
\begin{lstlisting}[language=Python]
for r, est in ((middle, []), (exit_, [(1, b"")])):
    st, hd = ntor3_client_start(r.id_ed, r.ntor, encode_extensions(est))
    body = extend2_body("127.0.0.1", r.port, r.id_rsa, r.id_ed,
                          HTYPE_NTOR_V3, hd)
    msg = RelayMessage(RelayCommand.EXTEND2, 0, body)
    c.send(Cell(CIRC, CellCommand.RELAY_EARLY, circ.crypto.encrypt(msg)))
    i, response = circ.crypto.decrypt(c.receive().body)
    sm, k = ntor3_client_finish(st, hdata_created2(response.data))
    circ.hops.append(Hop.from_keys(CircuitKeys.from_stream(k)))
\end{lstlisting}
Сообщение \texttt{EXTEND2} имеет StreamID 0, поскольку это управляющее сообщение, и передаётся в \cella{RELAY\_EARLY}. Для выходного реле мы также запрашиваем управление перегрузкой, как делает tor. Выполнение даёт:
\begin{lstlisting}[numbers=none]
переход 1: 003r
  <- CircID=0x80000100 RELAY тело=509 байт
  от перехода 1 EXTENDED2 67 байт
переход 2: 004r  расширения реле: []
  <- CircID=0x80000100 RELAY тело=509 байт
  от перехода 2 EXTENDED2 70 байт
переход 3: 006e  расширения реле: [(2, b'\x1f')]
\end{lstlisting}
Теперь цепь содержит три перехода. Каждая строка \enquote{от перехода} подтверждает, что наш код правильно зашифровал данные для нужного реле и что реле правильно зашифровало ответ. Ошибка хотя бы в одном ключе, байте дайджеста или порядке слоёв заставила бы реле уничтожить цепь. Длины \texttt{EXTENDED2} согласуются с \cref{sec:ntor3}: два байта HLEN, 32 для $Y$, 32 для $\mathit{AUTH}$ и сообщение $SM$, равное одному байту без расширений для промежуточного реле и четырём байтам с одним однобайтовым расширением для выходного реле.

\paragraph{Шаг 5: наконец заполнение.} Продолжаем слушать канал после расширения цепи:
\begin{lstlisting}[numbers=none]
  4.1 s  <- CircID=0x00000000 PADDING тело=509 байт
  7.1 s  <- CircID=0x00000000 PADDING тело=509 байт
 12.8 s  <- CircID=0x00000000 PADDING тело=509 байт
 21.5 s  <- CircID=0x00000000 PADDING тело=509 байт
 28.3 s  <- CircID=0x00000000 PADDING тело=509 байт
\end{lstlisting}
Именно это мы предсказали в главах \ref{cap:07-canali} и~\ref{cap:09-circuiti}: теперь, когда по каналу проходит многошаговая цепь, входное реле считает его \enquote{используемым} и при бездействии отправляет \cella{PADDING} через случайные интервалы. Наблюдавшиеся интервалы от трёх до 8,7 секунды укладываются в предусмотренное спецификацией окно 1,5--9,5 секунды.

% ---------------------------------------------------------------------------
\section*{Итоги}
\begin{itemize}
  \item В tor1 каждый переход имеет поток AES-128-CTR и накопительный дайджест SHA-1 для каждого направления; оба сохраняют состояние всю жизнь цепи.
  \item В прямом направлении клиент накладывает слои от последнего перехода к первому, а каждое реле снимает один. В обратном направлении каждое реле добавляет слой, а клиент снимает все.
  \item Клиент может адресовать любой переход (\emph{leaky pipe}).
  \item Ячейка распознаётся, когда Recognized равен нулю и накопительный дайджест совпадает; ложное срабатывание Recognized не должно менять состояние дайджеста.
  \item AES-CTR пластичен, а целостность проверяется только в пункте назначения: это делает возможным маркирование.
  \item CGO делает ячейки непластичными, связывает их теги, обновляет ключи после каждого сообщения и использует 128-битную аутентификацию; схема согласуется отдельно с каждым переходом посредством ntor-v3.
\end{itemize}

% ---------------------------------------------------------------------------
\section*{Упражнения}

\begin{esercizio}[1]\label{es:10-ordine}
Клиент с трёхпереходной цепью хочет отправить сообщение второму переходу. Какие ключи он использует для шифрования и в каком порядке? Что видит третий переход?
\end{esercizio}

\begin{esercizio}[1]\label{es:10-contatore}
После отправки клиентом 1000 ячеек по цепи в какой позиции своей гаммы шифрования AES-CTR находится первый переход? Почему тот же ключ нельзя использовать для второй цепи?
\end{esercizio}

\begin{esercizio}[2]\label{es:10-falso}
Объясните, что произошло бы, если реле обновляло накопительный дайджест даже тогда, когда Recognized равен нулю, но дайджест не совпадает. В среднем через сколько ячеек возникла бы проблема?
\end{esercizio}

\begin{esercizio}[2]\label{es:10-tag}
Сотрудничающее входное реле (guard) хочет пометить прямую ячейку так, чтобы сотрудничающее выходное реле её распознало. Какие байты ячейки было бы удобно изменить с учётом структуры открытого сообщения \texttt{DATA}? Что увидит честное выходное реле?
\end{esercizio}

\begin{esercizio}[2]\label{es:10-cgo-tag}
Повторите сценарий предыдущего упражнения с CGO. Что увидит сотрудничающее выходное реле? Что произойдёт с последующими ячейками цепи?
\end{esercizio}

\begin{esercizio}[2]\label{es:10-padding}
В шаге 5 лабораторной работы вычислите интервалы между наблюдавшимися ячейками \cella{PADDING}. Согласно спецификации, задержка выбирается из распределения максимума двух равномерных случайных величин. Какова ожидаемая задержка на интервале от 1,5 до 9,5 секунды? Согласуется ли она с наблюдениями?
\end{esercizio}

\begin{esercizio}[3]\label{es:10-drop}
Расширьте лабораторную работу, отправив сообщение \texttt{DROP} (команда RELAY 10, заполнение на всём протяжении цепи) второму переходу. Ответит ли реле? Как проверить, что оно распознало сообщение и не уничтожило цепь?
\end{esercizio}

\begin{esercizio}[3]\label{es:10-alterazione}
В лаборатории намеренно инвертируйте один бит зашифрованной ячейки \cella{RELAY\_EARLY}, содержащей \texttt{EXTEND2}. Какое реле обнаружит изменение и как отреагирует? Сначала предскажите результат, затем проверьте его.
\end{esercizio}

% ---------------------------------------------------------------------------
\section*{Дополнительная литература}
\begin{itemize}
  \item \titolo{Tor Protocol Specification}, главы \enquote{Relay cells} и \enquote{Routing relay cells}~\cite{torspec}.
  \item Предложение 359, \titolo{Counter Galois Onion, Updated}, со ссылками на работы Дегабриэле, Меллони, Мюнха и Стама.
  \item Предложения 202, 261, 295 и 308: многолетняя история подходов к предотвращению маркирования (tagging).
  \item R.~Dingledine, N.~Mathewson, P.~Syverson, \titolo{Tor: The Second-Generation Onion Router}~\cite{dingledine2004tor}, раздел 4 (цепи, leaky pipes и проверки целостности).
\end{itemize}
